% --- RESUMO EM PORTUGUÊS (264 palavras) ---
\begin{resumo}
Na Engenharia de Software, a Garantia da Qualidade (\textit{Quality Assurance} --- QA) e o gerenciamento estruturado de testes são essenciais para mitigar falhas em sistemas computacionais. Contudo, nos projetos de software do Instituto Federal do Espírito Santo (IFES) Campus Alegre, o controle de testes e o relato de erros ocorriam de forma descentralizada, utilizando documentos e planilhas compartilhadas em nuvem, o que ocasionava perda recorrente de histórico de execução e dificuldade de acompanhamento gerencial. Este trabalho teve como objetivo desenvolver e validar o \textbf{Test-Manager}, um sistema web centralizado para gerenciamento de casos de teste e rastreamento de defeitos. Metodologicamente, realizou-se uma pesquisa aplicada, guiada pela abordagem de Apresentação de Produto e pelo desenvolvimento iterativo e incremental, contemplando elicitação de requisitos, modelagem UML (casos de uso, classes e modelo lógico relacional) e implementação sob o padrão \textit{Model-View-Controller} (MVC) com PHP, Laravel, MySQL e Docker. A solução estrutura o fluxo de qualidade em três perfis de acesso (Coordenador, Tester e Desenvolvedor), abrangendo gestão de projetos, importação de roteiros em lote via CSV, versionamento automático de execuções (\texttt{v1.0}, \texttt{v2.0}), registro obrigatório de defeito com evidência visual na reprovação, trilha de auditoria e emissão de relatórios em PDF. A validação empírica, conduzida via questionário baseado no Modelo de Aceitação de Tecnologia (TAM) em escala Likert (1--5) aplicado a cinco participantes após uso orientado, indicou média de 4,43 (90,0\% de concordância) nas dimensões gerais, 4,80 (100,0\% de concordância) nos módulos específicos por perfil e média global de 4,53, apresentando indícios favoráveis de que a plataforma contribui para a organização, a rastreabilidade e a governança dos testes no ambiente acadêmico.

\textbf{Palavras-chave}: Engenharia de Software. Garantia da Qualidade. Teste de Software. Rastreamento de Defeitos. Laravel.
\end{resumo}

% --- ABSTRACT EM INGLÊS (225 palavras) ---
\begin{resumo}[Abstract]
 \begin{otherlanguage*}{english}
In Software Engineering, Quality Assurance (QA) and structured test management are essential activities to mitigate failures in computer systems. However, within software projects at the Federal Institute of Espírito Santo (IFES) Campus Alegre, test control and error reporting have historically relied on decentralized cloud documents and spreadsheets, leading to execution history loss and limited managerial visibility. This study aimed to develop and validate \textbf{Test-Manager}, a centralized web system for test case management and defect tracking. Methodologically, this applied research followed the Product Presentation approach and an iterative and incremental development process, encompassing requirements elicitation, UML modeling (use cases, class diagrams, and relational database design), and implementation under the Model-View-Controller (MVC) pattern using PHP, Laravel, MySQL, and Docker. The platform structures the QA workflow across three access profiles (Coordinator, Tester, and Developer), providing project management, batch test case import via CSV, automatic execution versioning (\texttt{v1.0}, \texttt{v2.0}), mandatory defect registration with visual evidence upon test rejection, an audit trail, and consolidated PDF reports. Empirical validation, conducted through a Technology Acceptance Model (TAM) questionnaire on a 5-point Likert scale with five participants after guided usage, yielded a mean score of 4.43 (90.0\% positive agreement) on general dimensions, 4.80 (100.0\% positive agreement) on role-specific modules, and an overall consolidated mean of 4.53, providing favorable evidence that the platform improves organization, historical traceability, and quality governance in academic software projects.

 \textbf{Keywords}: Software Engineering. Quality Assurance. Software Testing. Defect Tracking. Laravel.
 \end{otherlanguage*}
\end{resumo}